\documentclass[12pt]{article}

\usepackage{fontawesome}
\usepackage{hyperxmp}
\usepackage{hyperref}
\usepackage{xurl}
\usepackage[tagged, highstructure]{accessibility}
\usepackage{bookmark}
\usepackage{enumitem}

\hypersetup{
    pdfuapart=1,
    pdfdisplaydoctitle=true,
    colorlinks=true,
    urlcolor=[RGB]{0,40,150},
    linkcolor=[RGB]{0,40,150},
    citecolor=[RGB]{0,40,150},
    pdftitle={Basic Scripting},
    pdfauthor={Adrianna Holden-Gouveia},
    pdfsubject={Introduction to Linux - Lab Assignment},
    pdflang={en-US},
    pdfkeywords={lab assignment, Basic Scripting}
}

% Configure PDF bookmarks for navigation
\bookmarksetup{
    numbered,
    open,
}

% Configure list spacing for better accessibility
\setlist{nosep}



\title{Basic Scripting}
\author{
        Adrianna Holden-Gouveia \\
        Website: \href{https://aholdengouveia.name}{aholdengouveia.name}\\
        \date{\vspace{-5ex}}
        %Email: \href{mailto:admin@aholdengouveia.name}{admin@aholdengouveia.name} \\
%        \faLinkedin{: aholdengouveia} \\
        \faGithub {: aholdengouveia} \\
 %       \faTwitter {: aholdengouveia} \\
        }

%S\date{\today}


\begin{document}    

\maketitle

\section*{Accessibility Notice}
This document is also available in HTML format at:

Link: \href{https://aholdengouveia.name/IntroLinux/labs/shellscripting.html}{\textbf{https://aholdengouveia.name/IntroLinux/labs/shellscripting.html}}

The HTML version provides enhanced accessibility features including keyboard navigation, screen reader support, responsive design, dark mode support, and high contrast options.


%\begin{abstract}
%This is a template for Linux Administration Lab work
%\end{abstract}
%\tableofcontents

\section*{Objectives:}
\begin{itemize}
    \item Learn how to do some basic bash scripting. Become comfortable and competent with the steps needed to write, run, comment, and test scripts. 
\end{itemize}
\section*{Complete the following problems}

References, a video, a PowerPoint and some notes are available at my website.
Link: \href{https://www.aholdengouveia.name/IntroLinux/shellscripting.html}{\textbf{https://www.aholdengouveia.name/IntroLinux/shellscripting.html}}

\subsection*{How to turn in Scripts}
    \begin{enumerate}
        \item Testing must be submitted along with your code. You should submit a copy of the code and a screenshot or the results of your test so that I can see a successful run of your code. Every script should be tested.  A script that doesn't run is worth 0 points.
        \item Comments should be done on EVERY line of code, (this is very bad practice in the real world, never do this for a job) so that I can see that you understand the code that you are submitting. Every line should have an explanation of what that line does.  Comments are worth 5 points per script.
        \item For submission you need to submit a copy of your code (actual text NO screenshots of code allowed) and your successful test run of your code for every script.
\end{enumerate}

\subsection*{Scripts}
    \begin{enumerate}
        \item Write a script that will echo out the phrase “Hello World”  and then your name in plain text
        \item Write a script that will show each of the 3 types of echo (opaque, semi-opaque and cleartext (from the PowerPoint))\\
    \end{enumerate}

Make sure to make 2 separate scripts. Don't solve in just one.

\section*{Deliverables}
One file that includes all the requested information with your name, the date, and the lab title, make sure to follow the requested pattern for scripts. You need 2 separate scripts for this lab.

\end{document}